\documentclass[11pt]{article}

\usepackage[utf8]{inputenc}
\usepackage[T1]{fontenc}
\usepackage{amsmath,amssymb,amsfonts}
\usepackage{mathtools}
\usepackage{microtype}
\usepackage{graphicx}
\usepackage{hyperref}
\usepackage[margin=1in]{geometry}

\renewcommand{\topfraction}{0.9}
\renewcommand{\bottomfraction}{0.9}
\renewcommand{\textfraction}{0.1}
\renewcommand{\floatpagefraction}{0.8}
\renewcommand{\dbltopfraction}{0.9}
\renewcommand{\dblfloatpagefraction}{0.8}

\title{Set-Key Encoders for Discretization-Consistent Learning from Irregular Continuum Samples}

\author{
  Stepan Tretiakov \\
  University of Texas at Austin \\
  \texttt{stepan@utexas.edu}
}

\date{}

\emergencystretch=1em
\begin{document}
\maketitle

\begin{abstract}
Set-valued observations often arise as irregular samples of an underlying
continuum object.  In practice, both measurement budgets and training
computation are often limited.  Under these constraints, an encoder should
support training at a moderate sampling budget, improve systematically when
denser observations become available at test time without retraining, and
remain stable under changes in the sampling pattern.  We introduce
\emph{Set-Key}, a set encoder architecture
for irregular continuum samples.  In Set-Key,
coordinate-derived keys determine token--sensor interactions and value features
carry the observed signal, while the additive pooling stage defines token-wise
latent integral estimators.  This identifies latent-space integration as the
central architectural mechanism and exposes a clean insertion point for
geometry-aware discretization control.  We evaluate this idea under
same-object resampling protocols that vary both sample count and sampling
distribution while holding the underlying continuum signal fixed, using a
simple $k$-nearest-neighbor density correction as one concrete realization of
the geometry-aware pooling rule.  Across a synthetic surface-average
regression task, a real AhmedML surface-force benchmark, and a derived
AirfRANS airfoil-force regression benchmark, Set-Key improves both
matched-resolution robustness and zero-shot test-time refinement.  In all
three benchmarks, geometry-aware pooling reduces worst-case error under
sampling shift, lowers prediction drift, and makes better use of additional
test-time samples without retraining.  These results support latent-space
integration as a useful design principle for discretization-consistent set
encoding from irregular continuum samples.
\end{abstract}

\section{Introduction}
\label{sec:introduction}

Many machine learning problems receive set-valued observations that are, by
construction, finite samples of an underlying continuum object.  Examples
include sampled surface fields in computational fluid dynamics, irregular
boundary observations of aerodynamic profiles, and sensorized structural
surfaces.  In practice, training is often performed at a moderate sensor budget
because both observation and computational costs increase with input
resolution.  At
inference, however, the available number of sensors and the acquisition pattern
may differ from those seen during training.  A useful set encoder should
therefore train efficiently on sparse inputs, exploit denser test-time
observations without retraining, and remain reliable under shifts in the
sampling strategy.  Standard permutation-invariant encoders such as Deep
Sets~\cite{zaheer2017deep}, Set Transformer~\cite{lee2019set}, PointNet++~\cite{qi2017pointnetpp},
and PointNeXt~\cite{qian2022pointnext} aggregate over the empirical input set
and are typically evaluated under matched train/test sampling, so this
requirement is not built into their design.

The central difficulty is not only that the number of sensors may change, but
that the empirical sampling measure may change with it.  The same continuum
object may be observed through uniform, biased, clustered, or partially
occluded samples, and increasing the number of inference-time sensors does not
by itself remove this mismatch.  If aggregation remains tied to the observation
density, additional samples can reinforce the wrong latent summary rather than
improve it.  This issue is especially acute when the target depends on a
continuum quantity, such as a surface-averaged signal or a global force
coefficient, but it also affects any representation whose latent summaries
should depend primarily on the underlying object rather than on the particular
discretization.  We refer to this requirement as \emph{discretization
consistency}.

This paper addresses discretization consistency through a particular encoder
design.  We introduce the \emph{Set-Key} encoder, which decouples
coordinate-derived keys from value-derived signal features for irregular
continuum samples.  In Set-Key, token-wise latent summaries are
formed by an additive pooling stage, so each token can be interpreted as a
latent-space integral estimator of a learned geometry-conditioned test function
against the sampled signal.  This makes the pooling stage the natural location
for geometry-aware control of the sampling measure, while leaving the key,
value, and readout pathways unchanged.  Under this view, robustness to
sampling-strategy shift is not an auxiliary property, but a direct consequence
of how the encoder forms its latent summaries.

To assess discretization consistency directly, we introduce \emph{same-object resampling protocols}
that vary both sample count and sampling distribution while holding the
underlying continuum signal fixed.  The resulting evaluation targets the regime
of practical interest for this paper: models are trained once at a moderate
sensor budget and then tested under denser or differently distributed
observations of the same object, without retraining.  We report task accuracy
together with prediction drift, worst-case performance across sampling modes,
and refinement behaviour under increasing inference-time sample counts, so that
modelling quality and sensitivity to the empirical sampling measure can be
assessed separately.

We validate this perspective on three benchmarks spanning controlled and
real-data settings: synthetic surface-average regression on~$\mathbb{S}^2$,
AhmedML sparse surface-force regression, and AirfRANS sparse airfoil-force
regression.  Across all three, Set-Key with geometry-aware pooling is more
robust to sampling-strategy shift at matched resolution and benefits more
systematically from additional inference-time samples without retraining.  The
largest gains appear precisely under biased and nonuniform sampling regimes,
where standard additive aggregation remains most strongly tied to the empirical
sampling measure.

\section{Methods}
\label{sec:methods}

\subsection{Sampling Setup and Notation}
\label{sec:method-setup}

We treat each input set as a finite observation of an underlying continuum
surface or manifold signal.  Let $\mu$ denote the normalized target measure
over the object of interest, and let
$(\boldsymbol{x}_i,\boldsymbol{u}_i)_{i=1}^M$ be sampled coordinates and
values drawn from an observation process $\nu$ that may differ across views of
the same underlying object.  A model receives the sampled set
$S=\{(\boldsymbol{x}_i,\boldsymbol{u}_i)\}_{i=1}^M$ and predicts a task output
derived from the underlying continuum signal, such as a surface-average scalar
or a global force coefficient.  In the regime of interest for this paper, both
the sample count $M$ and the sampling pattern induced by $\nu$ may vary
between training and inference.

Our goal is not only high task accuracy under i.i.d.\ train/test sampling, but
stability under \emph{same-object resampling}: if two observation sets come
from the same underlying object but use different sample counts or density
profiles, their encoded representations and predictions should remain close.
This is the notion of discretization consistency studied throughout the paper.

\subsection{Set-Key Encoder Architecture}
\label{sec:method-architecture}

The core model studied in this paper is a permutation-invariant set encoder
that maps an unordered sensor set
$\{(\boldsymbol{x}_i,\boldsymbol{u}_i)\}_{i=1}^M$, with
$\boldsymbol{x}_i\in\mathbb{R}^{d_x}$ and
$\boldsymbol{u}_i\in\mathbb{R}^{d_u}$, to $p$ token-wise latent vectors
$\{\boldsymbol{b}_k\}_{k=1}^{p}$, where
$\boldsymbol{b}_k\in\mathbb{R}^{d_{\mathrm{out}}}$.  We refer to this structured encoder as a
\emph{Set-Key encoder}: token--sensor mixing is driven by a separate
position-only key pathway, while sensor content is carried by a separate value
pathway.  The key point for this paper is that these latent tokens are formed
by explicit additive aggregation, so they can be interpreted as learned
latent-space integrals of geometry-conditioned test functions against sampled
signals.

\paragraph{Input processing and sensor-side features.}
Let $\boldsymbol{e}_i$ denote the coordinate representation used by the
encoder.  In all experiments of this paper we take
$\boldsymbol{e}_i=\boldsymbol{x}_i$ directly.  A position-only Key Net
$k_{\boldsymbol{\theta}}:\mathbb{R}^{d_e}\to\mathbb{R}^{d_k}$ maps the encoded
sensor locations to key features
\[
\boldsymbol{K}_i = k_{\boldsymbol{\theta}}(\boldsymbol{e}_i)\in\mathbb{R}^{d_k},
\]
while a Value Net produces sensor-side value features
\[
\boldsymbol{V}_i = v_{\boldsymbol{\theta}}(\boldsymbol{u}_i)\in\mathbb{R}^{d_v}
\qquad\text{or}\qquad
\boldsymbol{V}_i = v_{\boldsymbol{\theta}}([\boldsymbol{e}_i,\boldsymbol{u}_i])\in\mathbb{R}^{d_v}.
\]
Both variants fit the present formulation.  In the experiments of this paper
we use the latter option, so the value pathway may depend on both geometry and
signal, while token--sensor mixing itself remains controlled solely by the
coordinate-derived keys.

\paragraph{Token--sensor aggregation.}
Let $\mathbf{Q}\in\mathbb{R}^{p\times d_k}$ denote a matrix of learnable query
tokens, and let
$\mathbf{K}=[\boldsymbol{K}_1,\dots,\boldsymbol{K}_M]^\top\in\mathbb{R}^{M\times d_k}$
and
$\mathbf{V}=[\boldsymbol{V}_1,\dots,\boldsymbol{V}_M]^\top\in\mathbb{R}^{M\times d_v}$
denote sensor-side key and value features, respectively.  The token--sensor
score matrix is
\[
\mathbf{S} = \frac{\mathbf{Q}\mathbf{K}^\top}{\sqrt{d_k}}
\in \mathbb{R}^{p \times M}.
\]
The implementation optionally applies a learned temperature and then a
pointwise basis nonlinearity $a(\cdot)$ to $\mathbf{S}$, with
$a\in\{\tanh,\mathrm{softsign},\mathrm{softplus}\}$.  In all experiments we use
the softplus choice and write
\[
\mathbf{B} = a(\mathbf{S}) \in \mathbb{R}^{p\times M}.
\]
Let
$\boldsymbol{w}\in\mathbb{R}_{+}^{M}$ denote sensor weights and
$\boldsymbol{m}\in\{0,1\}^{M}$ an optional validity mask.  Under the
``total'' normalization used throughout the experiments, the token matrix is
\begin{equation}
  \mathbf{P}
  =
  \frac{\mathbf{B}\,\operatorname{diag}(\boldsymbol{m}\odot\boldsymbol{w})\,\mathbf{V}}
       {\sum_{i=1}^{M} m_i w_i + \varepsilon}
  \in\mathbb{R}^{p\times d_v}.
  \label{eq:setkey_pool}
\end{equation}
Equivalently, for the $k$-th token,
\[
\boldsymbol{P}_k
\;=\;
\sum_{i=1}^{M} \tilde w_i\,a(S_{k,i})\,\boldsymbol{V}_i,
\qquad k=1,\dots,p.
\]
where
\[
\tilde w_i
\;=\;
\frac{m_i w_i}{\sum_{j=1}^{M} m_j w_j + \varepsilon},
\qquad i=1,\dots,M.
\]
The scalar weights $\boldsymbol{w}$ encode the geometry-aware correction of the
empirical sampling measure.  Uniform pooling is recovered by taking
$w_i \equiv 1$, while nonuniform choices of $\boldsymbol{w}$ alter only the
final additive pooling rule and are shared across all tokens.

\paragraph{Readout to token coefficients.}
The pooled token summaries are post-processed row-wise by a small map
$\rho_{\mathrm{tok}}:\mathbb{R}^{d_v}\to\mathbb{R}^{d_{\mathrm{out}}}$,
yielding
\[
\boldsymbol{b}_k = \rho_{\mathrm{tok}}(\boldsymbol{P}_k),
\qquad k=1,\dots,p.
\]
In the benchmarks considered here, the $p$ output tokens are flattened and
passed through a task head.

\subsection{Additive Pooling as a Latent-Space Integration Operator}
\label{sec:method-quadrature}

Equation~\eqref{eq:setkey_pool} identifies the additive pooling stage as the
site where the latent-integration interpretation becomes explicit.  Writing the
predictor as
\[
  \widehat{\boldsymbol{y}}(S)
  \;=\;
  r\!\bigl(\boldsymbol{P}_1(S),\dots,\boldsymbol{P}_p(S)\bigr),
\]
where $S=\{(\boldsymbol{x}_i,\boldsymbol{u}_i)\}_{i=1}^M$ is the observed set,
$\boldsymbol{P}_k(S)$ are the pooled token summaries from
Equation~\eqref{eq:setkey_pool}, and $r$ denotes the downstream readout
(combining the row-wise map $\rho_{\mathrm{tok}}$ with the final task head),
the quadrature-style interpretation applies to the intermediate token summaries
$\boldsymbol{P}_k$, and the readout acts on the resulting latent summaries.

For a fixed token $k$, the pooling rule has the form
\begin{equation}
  \boldsymbol{P}_k(S)
  \;=\;
  \sum_{i=1}^M \tilde w_i\,\phi_k(\boldsymbol{x}_i)\,\psi(\boldsymbol{x}_i,\boldsymbol{u}_i),
\end{equation}
where $\phi_k$ is a learned scalar test function indexed by token~$k$,
$\psi$ is a vector-valued value map, and $\tilde w_i$ are normalized
aggregation weights.  When the sampled set approximates a continuum signal
$\boldsymbol{u}:\mathcal{M}\to\mathbb{R}^{d_u}$ over a target measure $\mu$,
this expression is a quadrature-style estimator of the latent functional
\begin{equation}
  \bar{\boldsymbol{P}}_k
  \;=\;
  \int_{\mathcal{M}}
    \phi_k(\boldsymbol{x})\,
    \psi(\boldsymbol{x},\boldsymbol{u}(\boldsymbol{x}))\,
    \mathrm{d}\mu(\boldsymbol{x}).
\end{equation}
Accordingly, the readout $r$ operates on a finite collection of latent integral
estimators rather than on raw samples.

This decomposition yields an immediate stability relation.  If the readout
$r:\mathbb{R}^{p d_v}\to\mathbb{R}^{d_{\mathrm{out}}}$ is
$L_r$-Lipschitz under the Euclidean norm, and if
\[
  \mathbf{P}(S)
  =
  \operatorname{vec}\bigl[\boldsymbol{P}_1(S),\dots,\boldsymbol{P}_p(S)\bigr]
  \in \mathbb{R}^{p d_v},
\]
then token-level integration error controls output error:
\[
  \|\widehat{\boldsymbol{y}}(S)-\widehat{\boldsymbol{y}}(S')\|
  \;\le\;
  L_r\,
  \bigl\|
    \mathbf{P}(S)-\mathbf{P}(S')
  \bigr\|.
\]
Hence, improvements in the additive pooling rule improve the quality of the
latent integral estimators supplied to the readout, and the induced stability
gain propagates through the final predictor.

Uniform pooling corresponds to the special case $\tilde w_i=1/M$, which is
correct only when the observation process already matches the target measure
$\mu$.  If the same object is observed under a shifted sampling density
$\nu \neq \mu$, uniform pooling converges to the wrong latent functional even
though permutation invariance is preserved.  This token-level mismatch is the
basic source of discretization-induced prediction drift.

\subsection{Geometry-Aware Weighting via kNN Density Estimates}
\label{sec:method-aqp}

As a simple geometry-aware baseline, we estimate local sampling density from
$k$-nearest-neighbor radii and use those estimates to reweight the additive
pooling rule.  If $r_k(\boldsymbol{x}_i)$ denotes the distance from point
$\boldsymbol{x}_i$ to its
$k$-th nearest neighbor, then for an intrinsic dimension $d$ we define
\begin{equation}
  w_i \;\propto\; r_k(\boldsymbol{x}_i)^d,
\end{equation}
followed by normalization so that $\sum_i w_i = 1$ over the observed set.
Intuitively, densely sampled regions have smaller $r_k$ and therefore receive
less mass, while sparse regions receive more.  This makes the rule a crude
local cell-area estimator and therefore a natural baseline when the target
measure differs from the observation density.

We do not regard this $k$-NN rule as a final answer.  It is a useful
heuristic density correction and an effective empirical baseline, but it
captures only one scalar notion of local scale and does not explicitly model
anisotropy, curvature, or missing support.  One goal of the benchmark suite is
to quantify exactly how far such simple geometry-aware weighting goes and where
stronger aggregation rules will be needed.

\subsection{Discretization Consistency}
\label{sec:method-consistency}

We call an encoder discretization consistent if its output depends primarily on
the underlying continuum object and only weakly on the specific finite
discretization used to observe that object.  Formally, if two sets
$S,S'$ are sampled from the same continuum object under different resolutions
or densities, then a discretization-consistent encoder should satisfy
$E(S)\approx E(S')$ and therefore produce stable downstream predictions.

This requirement is distinct from ordinary generalization.  A model can achieve
high i.i.d.\ accuracy while still reacting strongly to harmless changes in
sampling density.  Our empirical metrics therefore separate task quality from
representation stability: prediction drift measures output sensitivity under
same-object resampling, and
worst-case performance across sampling modes measures robustness.

For a test object $o$, let $S_o^{\mathrm{ref}}$ denote a dense uniform
reference view and let $S_{o,m,N,r}$ denote a resampled view under sampling
mode $m$, point count $N$, and replica $r$.  Writing
$\widehat{\boldsymbol{y}}_o^{\mathrm{ref}}=\widehat{\boldsymbol{y}}(S_o^{\mathrm{ref}})$
and
$\widehat{\boldsymbol{y}}_{o,m,N,r}=\widehat{\boldsymbol{y}}(S_{o,m,N,r})$,
we define the prediction drift for that setting as
\[
  D_{m,N}
  \;=\;
  \frac{1}{|T|\,R\,d_{\mathrm{out}}}
  \sum_{o\in T}
  \sum_{r=1}^{R}
  \bigl\|
    \widehat{\boldsymbol{y}}_{o,m,N,r}
    -
    \widehat{\boldsymbol{y}}_o^{\mathrm{ref}}
  \bigr\|_1,
\]
where $T$ is the test set, $R$ is the number of resampling replicas, and
$d_{\mathrm{out}}$ is the output dimension.  Thus, the table entry labelled
``Pred.\ Drift'' is the mean absolute change in prediction relative to the
dense reference view of the same object.

\section{Experiments}
\label{sec:experiments}

The experiments evaluate Set-Key in the regime that motivates this paper:
training at a moderate sensor budget, followed by inference under denser or
differently distributed observations of the same underlying object without
retraining.  The benchmark suite covers three continuum settings: a controlled
synthetic surface-functional regression task on~$\mathbb{S}^2$, a real AhmedML
surface-force regression benchmark, and a derived AirfRANS airfoil-force
regression benchmark.  Within the Set-Key family we compare two matched models
that differ only in the pooling weights: a \textbf{uniform} encoder with
standard additive pooling and a \textbf{$k$-NN density} encoder with
geometry-aware weighting ($k{=}8$).  All models are trained on uniformly
sampled observations and then evaluated under same-object resampling protocols
that vary both the sampling pattern and the number of observed samples.  We
report task quality together with \emph{discretization-consistency}
metrics---prediction agreement with a high-resolution reference, prediction
drift, and worst-case performance across sampling modes---so that modelling
quality and stability under sampling shift are assessed separately.

For architecture-level context we also report results with a
PointNeXt-style~\cite{qian2022pointnext} backbone as a strong reference model
with learned local aggregation.  This comparison tests whether the gains from
geometry-aware pooling in Set-Key remain visible relative to a competitive
modern point-set architecture, rather than only relative to uniform variants of
the same encoder.  Because PointNeXt aggregates features through learned local
neighborhood blocks and does not expose an explicit scalar-weight insertion
point analogous to the $\boldsymbol{w}$ in Equation~\eqref{eq:setkey_pool}, we
treat it as a reference architecture rather than as a geometry-weighted variant
of Set-Key.

\paragraph{Sampling strategies.}
Across all benchmarks, a sampling strategy is a named resampling distribution
obtained by reweighting the benchmark's base observation measure while holding
the underlying continuum object and target fixed.  The base measure is uniform
surface measure on~$\mathbb{S}^2$ for the synthetic benchmark, mesh-cell area
on AhmedML, and boundary arclength on AirfRANS.  We organize these strategies
into four cross-benchmark families.  First, \emph{uniform} leaves the base
measure unchanged.  Second, \emph{region-biased} strategies upweight a named
part of the object: on~$\mathbb{S}^2$ this appears as \emph{polar} and
\emph{equatorial} bias, on AhmedML as \emph{front}, \emph{rear}, and
\emph{roof}, and on AirfRANS as \emph{leading edge}, \emph{trailing edge},
\emph{upper surface}, and \emph{lower surface}.  Third, \emph{clustered}
strategies concentrate samples around one or two random local centers; this is
the mechanism used by the shared \emph{clustered} mode on all benchmarks.
Fourth, \emph{partial-view} strategies suppress part of the support while
leaving the underlying object fixed: on~$\mathbb{S}^2$ this is realized by the
\emph{hemisphere} mode induced by a random viewing direction, while on AhmedML
and AirfRANS it is realized by the shared \emph{occluded} mode, which
downweights a random local neighborhood.  We therefore reuse names only when
the underlying mechanism matches across benchmarks, most directly for
\emph{uniform}, \emph{clustered}, and \emph{occluded}; benchmark-specific
names are reserved for domain-specific region-biased views.  Later subsections
refer to both families and concrete strategy names while keeping the
evaluation protocol fixed: for each strategy we evaluate multiple point counts
and independent replicas, and compare each prediction against a dense uniform
reference view of the same object.

\paragraph{Quantitative summaries.}
Tables~\ref{tab:main-results}, \ref{tab:sampling-breakdown},
\ref{tab:resolution-refinement}, \ref{tab:airfrans-sampling-breakdown}, and
\ref{tab:airfrans-resolution-refinement} provide the main quantitative
summaries.  Table~\ref{tab:main-results} gives the cross-benchmark comparison
for Set-Key across synthetic surface-functional regression together with the
AhmedML and AirfRANS real-data force-regression benchmarks.  Tables
\ref{tab:sampling-breakdown} and \ref{tab:resolution-refinement} summarize the
matched-resolution sampling breakdown and refinement behaviour on AhmedML,
while Tables~\ref{tab:airfrans-sampling-breakdown} and
\ref{tab:airfrans-resolution-refinement} provide the corresponding summaries
for AirfRANS.

% Auto-generated by paper/generate_result_tables.py
\begin{table*}[t]
\centering
\small
\setlength{\tabcolsep}{4.0pt}
\resizebox{\textwidth}{!}{%
\begin{tabular}{lccccccccc}
\hline
 & \multicolumn{3}{c}{\shortstack{Synthetic\\Regr.}} & \multicolumn{3}{c}{AhmedML} & \multicolumn{3}{c}{AirfRANS} \\
Model & \shortstack{In-dist.\\RMSE} & \shortstack{Worst\\RMSE} & \shortstack{Pred.\\Drift} & \shortstack{In-dist.\\RMSE} & \shortstack{Worst\\RMSE} & \shortstack{Pred.\\Drift} & \shortstack{In-dist.\\RMSE} & \shortstack{Worst\\RMSE} & \shortstack{Pred.\\Drift} \\
\hline
Set-Key (Unif) & 0.046 & 0.248 & 0.159 & 0.023 & 0.065 & 0.016 & 0.022 & 0.100 & 0.036 \\
Set-Key (kNN) & 0.021 & 0.188 & 0.087 & \textbf{0.022} & 0.037 & \textbf{0.008} & \textbf{0.017} & \textbf{0.056} & \textbf{0.013} \\
Set-Value (Vor) & \textbf{0.012} & \textbf{0.137} & \textbf{0.059} & -- & -- & -- & -- & -- & -- \\
PointNeXt & 0.042 & 0.146 & 0.171 & 0.024 & \textbf{0.037} & 0.031 & 0.024 & 0.063 & 0.219 \\
\hline
\end{tabular}
}
\caption{Main matched-resolution results. Each task is evaluated at its training resolution, so this table isolates density-shift robustness from resolution shift. We report in-distribution RMSE, worst-case RMSE, and average nonuniform prediction drift for both the synthetic surface-functional task together with the two real-data benchmarks: AhmedML sparse surface-force regression and AirfRANS sparse airfoil-force regression. Here ``In-dist.\ RMSE'' denotes RMSE under uniform sampling at the training resolution, ``Worst RMSE'' denotes the largest RMSE across sampling modes at that same resolution, and ``Pred.\ Drift'' denotes the mean absolute difference between a prediction from a sparse or shifted view and the dense-reference prediction for the same object, averaged over target components and over nonuniform sampling modes. The Set-Value (Vor) row is reported only for the synthetic sphere benchmark, where exact spherical Voronoi cell areas under the target measure are available.}
\label{tab:main-results}
\end{table*}


\subsection{Surface-Functional Regression under Resampling}
\label{sec:exp-regression}

Our first controlled benchmark targets a continuum functional with a known
measure: the surface average of a scalar field on~$\mathbb{S}^2$.  This setting
isolates the role of aggregation weights particularly clearly, since the target
depends directly on integration over the underlying surface rather than on a
downstream semantic label.

\paragraph{Setup.}
Each object is a continuous scalar field on the unit sphere~$\mathbb{S}^2$ of
the form
\begin{equation}\label{eq:sphere-field}
  f(\mathbf{p})
  \;=\;
  b
  + \mathbf{p}^\top \boldsymbol{\ell}
  + \boldsymbol{\phi}(\mathbf{p})^\top \mathbf{q}
  + \sum_{j=1}^{4}
      a_j \exp\!\bigl[\kappa_j\bigl(\mathbf{p}^\top \mathbf{c}_j - 1\bigr)\bigr],
\end{equation}
where $b$ is a scalar bias, $\boldsymbol{\ell}\in\mathbb{R}^3$ and
$\mathbf{q}\in\mathbb{R}^5$ are linear and quadratic coefficients over
hand-crafted polynomial features~$\boldsymbol{\phi}$, and the final sum adds
four localized Gaussian-like bumps with random centers
$\mathbf{c}_j\in\mathbb{S}^2$, amplitudes~$a_j$, and
concentrations~$\kappa_j\in[8,22]$.  All parameters are drawn independently
for each object.  We generate 2\,048 training, 256 validation, and 512 test
objects.  The target quantity is the surface average
$\bar{f}=\int_{\mathbb{S}^2} f\,\mathrm{d}\sigma$, approximated using 4\,096
uniform reference points.

Both regressors share the same weighted set encoder architecture: 16 learnable
query tokens, key dimension 64, token dimension 32, hidden dimension 128, and
a three-layer MLP head.  They differ only in their aggregation
weights---$1/N$ or $k$-NN density estimates ($k{=}8$, intrinsic
dimension~$d{=}2$).  Each model is trained for 2\,000 steps on batches of 64
objects, each observed through 128 uniformly sampled points, using AdamW
(learning rate~$10^{-3}$, weight decay~$10^{-5}$) with gradient clipping at
norm~1.

\paragraph{Evaluation protocol.}
At test time we resample each object under the five sphere sampling strategies
defined above: one uniform strategy, two region-biased strategies
(\emph{polar} and \emph{equatorial}), one clustered strategy
(\emph{clustered}), and one partial-view strategy (\emph{hemisphere}), at each
of five point counts
$N\in\{32,64,128,256,512\}$, with three independent resampling replicas per
setting.  For every object we also obtain a \emph{reference prediction} from a
dense 2\,048-point uniform sample.  This reference is not ground truth; it is
the model's own best-informed output, and serves as an anchor for measuring how
much the prediction shifts when the input is degraded or redistributed.

\paragraph{Metrics.}
We report four quantities per setting:
\emph{RMSE} (root mean squared error against the Monte Carlo estimate of the surface average),
\emph{MAE} (mean absolute error),
\emph{bias} (signed mean error, indicating systematic over- or
under-prediction), and
\emph{prediction drift} (mean absolute difference from the reference
prediction).
The first three characterize task-level accuracy; the last isolates
discretization instability.  A model can achieve low bias yet high drift---it
is correct on average but gives a different answer each time the same object
is resampled.  This error--drift decomposition is the central diagnostic:
it separates estimation quality from representation stability in a way
that scalar accuracy alone cannot.  We summarize with worst-case RMSE across
modes and average non-uniform prediction drift.

\paragraph{Observed pattern.}
This synthetic regression task is the controlled case where the quadrature
interpretation is most directly aligned with the target: the model is asked to
recover a surface average from irregular samples.  Table~\ref{tab:main-results}
therefore gives a matched-resolution summary.  At the training resolution,
$k$-NN weighting lowers in-distribution RMSE from $0.046$ to $0.021$, lowers
worst-case RMSE under sampling shift from $0.248$ to $0.188$, and reduces
average nonuniform prediction drift from $0.159$ to $0.087$.  The same trend is
visible under refinement in Figure~\ref{fig:regression-convergence}: the
geometry-aware encoder follows a consistently lower error curve and exhibits a
steeper reduction in same-object prediction drift as the number of observed
points grows, while PointNeXt serves as a larger-backbone reference rather than
an explicitly weighted quadrature model.
Figure~\ref{fig:regression-qualitative} shows the same comparison at the
prediction level for one challenging density shift: under polar sampling at
$N{=}512$, the $k$-NN encoder produces a noticeably tighter true-versus-predicted
scatter and a smaller residual cloud than uniform pooling, while the PointNeXt
reference remains competitive but more dispersed on this synthetic functional
estimation task.

\begin{figure*}[ht]
  \centering
  \includegraphics[width=\textwidth]{../results/point_cloud_mean_regression_run/point_cloud_mean_regression_convergence.pdf}
  \caption{Synthetic mean-regression refinement curves. Average nonuniform RMSE (left), shown here as a smoother refinement diagnostic complementary to the worst-case table metric, and average nonuniform prediction drift (right), defined as the mean absolute change from the dense-reference prediction of the same object, as the number of observed sphere samples increases. Dotted lines show log-log power-law fits.}
  \label{fig:regression-convergence}
\end{figure*}

\begin{figure*}[ht]
  \centering
  \includegraphics[width=\textwidth]{../results/point_cloud_mean_regression_run/point_cloud_mean_regression_prediction.pdf}
  \caption{Synthetic mean-regression predictions under polar sampling at $N{=}512$. A representative test sphere colored by the ground-truth scalar field with observed samples overlaid (left), true-versus-predicted scatter for the Monte Carlo surface-average target (middle), and residuals relative to that target (right) for uniform pooling, $k$-NN weighting, and PointNeXt.}
  \label{fig:regression-qualitative}
\end{figure*}

\subsection{AhmedML Surface-Force Regression}
\label{sec:exp-ahmed}

\paragraph{Task.}
To test whether the same aggregation mechanism transfers beyond the synthetic
sphere setting, we consider a real sparse surface-functional benchmark derived
from the AhmedML dataset of Ahmed-body CFD simulations.  Each example consists
of a triangulated vehicle surface together with the corresponding force
coefficients.  We
preprocess each mesh into an irregular set of cell-centered surface samples
with associated area weights, using cell centers as coordinates and
concatenating local pressure, cell normals, and wall-shear vectors as the
observed per-cell signal.  The targets are the drag and lift coefficients
$(C_d, C_l)$, so the task is multivariate regression of physically meaningful
global surface functionals from sparse irregular surface observations.

\paragraph{Setup and resampling protocol.}
The benchmark uses all 500 processed AhmedML runs, split into 350 training, 75
validation, and 75 test bodies.  Each model is trained on 256 uniformly
sampled surface cells per object and evaluated under the six AhmedML sampling
strategies defined above: one uniform strategy, three region-biased strategies
(\emph{front}, \emph{rear}, and \emph{roof}), one clustered strategy
(\emph{clustered}), and one partial-view strategy (\emph{occluded}).  The
nonuniform strategies are obtained by reweighting the base surface-area
distribution, so the sampling pattern changes while the underlying CFD
solution remains fixed.  We evaluate at point counts
$N\in\{64,128,256,512,1024\}$ with three independent resampling replicas per
setting, and use a dense 2\,048-cell reference view to measure same-object
prediction drift.

\paragraph{Metrics.}
We summarize AhmedML with the same measure-sensitive diagnostics used in the
synthetic regression task: in-distribution RMSE at the matched training
resolution, worst-case RMSE across sampling modes, and average nonuniform
prediction drift from the dense reference prediction.  Table~\ref{tab:main-results}
reports the cross-task summary, Table~\ref{tab:sampling-breakdown} gives the
per-mode breakdown at the matched training resolution $N{=}256$, and
Table~\ref{tab:resolution-refinement} shows how worst-case error and drift
change as the sampling budget grows.

% Auto-generated by paper/generate_result_tables.py
\begin{table*}[t]
\centering
\small
\setlength{\tabcolsep}{4.2pt}
\begin{tabular}{lcccccccc}
\hline
Model & uniform & front & rear & roof & clustered & occluded & Avg. Nonunif. & Worst \\
\hline
Set-Key (Unif) & 0.023 & 0.065 & 0.029 & 0.028 & 0.025 & 0.023 & 0.034 & 0.065 \\
Set-Key (kNN) & \textbf{0.022} & 0.037 & \textbf{0.021} & \textbf{0.023} & \textbf{0.023} & \textbf{0.021} & \textbf{0.025} & 0.037 \\
PointNeXt & 0.024 & \textbf{0.036} & 0.026 & 0.037 & 0.025 & 0.025 & 0.030 & \textbf{0.037} \\
\hline
\end{tabular}
\caption{Sampling-mode breakdown on the AhmedML benchmark at the matched training resolution $N=256$. Entries report RMSE under each sampling mode, followed by the average over nonuniform modes and the worst-case value across all modes.}
\label{tab:sampling-breakdown}
\end{table*}

% Auto-generated by paper/generate_result_tables.py
\begin{table*}[t]
\centering
\small
\setlength{\tabcolsep}{4.0pt}
\begin{tabular}{lcccccccccc}
\hline
 & \multicolumn{5}{c}{Worst RMSE $\downarrow$} & \multicolumn{5}{c}{Pred. Drift $\downarrow$} \\
Model & 64 & 128 & 256 & 512 & 1024 & 64 & 128 & 256 & 512 & 1024 \\
\hline
Set-Key (Unif) & 0.080 & 0.069 & 0.065 & 0.053 & 0.054 & 0.023 & 0.018 & 0.016 & 0.014 & 0.013 \\
Set-Key (kNN) & \textbf{0.075} & \textbf{0.045} & 0.037 & \textbf{0.023} & \textbf{0.024} & \textbf{0.021} & \textbf{0.013} & \textbf{0.008} & \textbf{0.005} & \textbf{0.003} \\
PointNeXt & 0.119 & 0.069 & \textbf{0.037} & 0.033 & 0.061 & 0.052 & 0.032 & 0.031 & 0.027 & 0.017 \\
\hline
\end{tabular}
\caption{Resolution/refinement study on the AhmedML benchmark. The left block reports the primary worst-case task metric across point counts, while the right block reports the corresponding same-object consistency metric, namely average nonuniform prediction drift relative to the dense-reference prediction for the same object. This table is intended to expose whether gains from geometry-aware weighting persist as the discretization is refined.}
\label{tab:resolution-refinement}
\end{table*}


\paragraph{Observed pattern.}
AhmedML is the primary real-data case in the paper where geometry-aware
pooling helps on the task metric itself rather than only on robustness
proxies.  At the matched training resolution $N{=}256$, $k$-NN weighting
reduces worst-case RMSE from approximately $0.065$ to $0.037$ and average
nonuniform prediction drift from approximately $0.016$ to $0.008$, with the
largest gains appearing under the front and clustered sampling modes.
Figure~\ref{fig:ahmed-convergence}
shows the same effect as a refinement plot: the geometry-aware encoder follows
a steeper decay curve for both worst-case error and same-object
prediction drift as the surface-cell budget grows, even against the PointNeXt
reference architecture.

\begin{figure*}[ht]
  \centering
  \includegraphics[width=\textwidth]{../results/ahmedml_surface_forces_run/ahmedml_surface_forces_convergence.pdf}
  \caption{AhmedML refinement curves. Worst-case RMSE across sampling modes (left) and average nonuniform prediction drift (right), defined as the mean absolute change from the dense-reference prediction of the same object, as the surface-cell budget grows. Dotted lines show log-log power-law fits.}
  \label{fig:ahmed-convergence}
\end{figure*}

To complement the aggregate tables, we also include one qualitative prediction
figure for AhmedML.  Figure~\ref{fig:ahmed-qualitative} combines a
representative front-biased surface view with test-set scatter plots for
$C_d$ and $C_l$ under the same sampling condition at the matched training
resolution $N{=}256$.  The figure is shown at a denser inference condition
$N{=}1024$ to illustrate the same-object prediction quality under refinement.
This makes the benchmark visually concrete while also showing the actual
regression quality of the two Set-Key variants together with a PointNeXt
reference.

\begin{figure*}[ht]
  \centering
  \includegraphics[width=\textwidth]{../results/ahmedml_surface_forces_run/ahmedml_surface_forces_prediction.pdf}
  \caption{AhmedML predictions under front-biased sampling at $N{=}1024$. A representative test body with observed cells overlaid (left), and true-vs-predicted scatter for $C_d$ (middle) and $C_l$ (right) comparing the uniform Set-Key variant, the geometry-aware $k$-NN Set-Key variant, and PointNeXt.}
  \label{fig:ahmed-qualitative}
\end{figure*}

\subsection{AirfRANS Airfoil-Force Regression}
\label{sec:exp-airfrans}

\paragraph{Task.}
To complement AhmedML with a second real aerodynamic functional benchmark, we
derive a sparse airfoil-force regression task from the raw AirfRANS CFD
simulations.  Each example is an irregular sample of the airfoil boundary
curve.  We preprocess every simulation into an airfoil point set with
coordinates, pressure, wall-shear stress, and boundary normals, and take the
global targets to be the aerodynamic force coefficients $(C_d, C_l)$ computed
with the AirfRANS reference API.

\paragraph{Setup and resampling protocol.}
We use the official AirfRANS ``full'' split as prepared by the benchmark
preprocessing pipeline, yielding 720 training, 80 validation, and 200 test
simulations.  Each model is trained for 25\,000 steps on 256 uniformly sampled
airfoil boundary points and then evaluated under the seven AirfRANS sampling
strategies defined above: one uniform strategy, four region-biased strategies
(\emph{leading edge}, \emph{trailing edge}, \emph{upper surface}, and
\emph{lower surface}), one clustered strategy (\emph{clustered}), and one
partial-view strategy (\emph{occluded}).  These strategies reweight the base
boundary-length distribution, so the observed subset shifts around the same
underlying airfoil solution while the reference coefficients remain fixed.
We test at point counts
$N\in\{64,128,256,512,1024\}$ and use a dense 2\,048-point reference view for
same-simulation prediction drift.

\paragraph{Metrics.}
For AirfRANS we report the same diagnostics used for AhmedML:
matched-resolution RMSE on $(C_d, C_l)$ at $N{=}256$, worst-case RMSE across
sampling modes, and average nonuniform prediction drift from the dense
reference prediction on the same simulation.  Table~\ref{tab:main-results}
includes the matched-resolution summary; Table~\ref{tab:airfrans-sampling-breakdown}
and Table~\ref{tab:airfrans-resolution-refinement} give the detailed
sampling-mode and refinement breakdowns.

% Auto-generated by paper/generate_result_tables.py
\begin{table*}[t]
\centering
\small
\setlength{\tabcolsep}{4.2pt}
\begin{tabular}{lccccccccc}
\hline
Model & uniform & \shortstack{leading\\edge} & \shortstack{trailing\\edge} & \shortstack{upper\\surface} & \shortstack{lower\\surface} & clustered & occluded & Avg. Nonunif. & Worst \\
\hline
Uniform & 0.022 & 0.081 & 0.100 & 0.078 & 0.056 & 0.042 & 0.024 & 0.063 & 0.100 \\
kNN & \textbf{0.017} & \textbf{0.028} & \textbf{0.029} & 0.056 & 0.043 & \textbf{0.022} & \textbf{0.021} & \textbf{0.033} & \textbf{0.056} \\
PointNeXt & 0.024 & 0.058 & 0.063 & \textbf{0.028} & \textbf{0.032} & 0.031 & 0.027 & 0.040 & 0.063 \\
\hline
\end{tabular}
\caption{Sampling-mode breakdown on the AirfRANS benchmark at the matched training resolution $N=256$. Entries report RMSE under each sampling mode, followed by the average over nonuniform modes and the worst-case value across all modes.}
\label{tab:airfrans-sampling-breakdown}
\end{table*}

% Auto-generated by paper/generate_result_tables.py
\begin{table*}[t]
\centering
\small
\setlength{\tabcolsep}{4.0pt}
\begin{tabular}{lcccccccccc}
\hline
 & \multicolumn{5}{c}{Worst RMSE $\downarrow$} & \multicolumn{5}{c}{Pred. Drift $\downarrow$} \\
Model & 64 & 128 & 256 & 512 & 1024 & 64 & 128 & 256 & 512 & 1024 \\
\hline
Uniform & \textbf{0.146} & \textbf{0.120} & 0.100 & 0.077 & 0.157 & \textbf{0.050} & 0.044 & 0.036 & 0.023 & 0.064 \\
kNN & 0.193 & 0.134 & \textbf{0.056} & \textbf{0.018} & \textbf{0.019} & 0.054 & \textbf{0.030} & \textbf{0.013} & \textbf{0.005} & \textbf{0.006} \\
PointNeXt & 0.198 & 0.136 & 0.063 & 0.129 & 0.349 & 0.206 & 0.205 & 0.219 & 0.192 & 0.094 \\
\hline
\end{tabular}
\caption{Resolution/refinement study on the AirfRANS benchmark. The left block reports the primary worst-case task metric across point counts, while the right block reports the corresponding same-object consistency metric, namely average nonuniform prediction drift relative to the dense-reference prediction for the same object. This table is intended to expose whether gains from geometry-aware weighting persist as the discretization is refined.}
\label{tab:airfrans-resolution-refinement}
\end{table*}


\paragraph{Observed pattern.}
This derived AirfRANS force-regression task supports the same matched-resolution
conclusion as AhmedML, while showing a more heterogeneous refinement pattern at
very low sampling budgets.  At the matched training resolution
$N{=}256$, $k$-NN weighting lowers in-distribution RMSE from approximately
$0.0218$ to $0.0168$, lowers worst-case RMSE from approximately $0.0998$ to
$0.0565$, and reduces average nonuniform prediction drift from approximately
$0.0359$ to $0.0130$.  The PointNeXt reference is competitive on worst-case
RMSE at the matched resolution, but it remains much less discretization-stable,
with drift above $0.21$.  Under refinement, the $k$-NN encoder reaches
worst-case RMSE below $0.019$ and prediction drift below $0.006$ by
$N{=}1024$, while both uniform pooling and PointNeXt deteriorate under the most
adverse nonuniform views.  Figure~\ref{fig:airfrans-convergence} shows the
same empirical refinement trend directly.  Because the reported
quantity is the worst-case error across several resampling modes at each point
count, monotonic decay is not expected for every model; the relevant feature is
the separation between methods as the sampling budget increases.  In that
comparison, the geometry-aware encoder separates from both uniform pooling and
PointNeXt in both worst-case error and same-object drift.  Figure~\ref{fig:airfrans-prediction}
shows the same comparison at the per-simulation level under a
trailing-edge sampling pattern: the left panel visualizes a representative test
airfoil with the observed boundary samples overlaid, while the scatter plots on
the right show that $k$-NN weighting stays much closer to the true $(C_d,C_l)$
targets than either uniform pooling or the PointNeXt reference.

\begin{figure*}[ht]
  \centering
  \includegraphics[width=\textwidth]{../results/airfrans_field_prediction_run/airfrans_force_regression_convergence.png}
  \caption{AirfRANS refinement curves. Worst-case RMSE across sampling modes (left) and average nonuniform prediction drift (right), defined as the mean absolute change from the dense-reference prediction of the same simulation, as the sampled airfoil-point budget grows. Dotted lines show log-log power-law fits.}
  \label{fig:airfrans-convergence}
\end{figure*}

\begin{figure*}[ht]
  \centering
  \includegraphics[width=\textwidth]{../results/airfrans_field_prediction_run/airfrans_force_regression_prediction.png}
  \caption{AirfRANS predictions under trailing-edge-biased sampling at $N{=}1024$. Left: a representative test airfoil colored by boundary pressure, with the observed boundary samples overlaid. Middle and right: true-versus-predicted scatter plots for $C_d$ and $C_l$ across the full AirfRANS test split, comparing uniform pooling, the geometry-aware $k$-NN encoder, and PointNeXt under the same observation pattern. The diagonal line indicates perfect prediction.}
  \label{fig:airfrans-prediction}
\end{figure*}

\clearpage
\bibliographystyle{plain}
\bibliography{references}

\end{document}
